\documentclass[9pt,twocolumn]{article}
\usepackage[a4paper,margin=16mm,columnsep=5mm]{geometry}
\usepackage{booktabs}
\usepackage{graphicx}
\usepackage[hidelinks]{hyperref}
\usepackage{xcolor}

\newcommand{\modeltag}{\textsuperscript{[model]}}
\newcommand{\pending}{\textcolor{red!70!black}{[PENDING]}}
\newcommand{\boundary}[1]{\textit{Boundary: #1}}
\newcommand{\placeholder}[1]{%
  \begin{center}\fbox{\parbox{0.90\columnwidth}{\centering\vspace{5mm}
  \textbf{Claim-safe placeholder}\par #1\vspace{5mm}}}\end{center}}

% Claim-safe evidence table.  Update only from sealed authorities.
\newcommand{\EvidenceRegistry}{%
\begin{center}
\scriptsize
\begin{tabular}{p{0.14\columnwidth}p{0.34\columnwidth}p{0.42\columnwidth}}
\toprule
Item & Admitted value & Mandatory boundary \\
\midrule
C1 ledger & 1.7591725402$\times$; 434,242,823 vs 763,908,050 &
CPU same-ledger component model; ten samples; not RTL/mapped/system \\
C1 physical & 147,246.39209~$\mu$m$^2$; setup +0.001795~ns &
28 nm, 3 ns, nine-macro slice; hold/full storage/power open \\
C2 & 1.01672765$\times$ cycles; 4.541077998$\times$ throughput/mm$^2$;
$-77.6104\%$ logic area &
K8 vs equal-bandwidth K1$\times$8; five directed workloads; logic-only pre-macro \\
C3 & 62,433.503388~$\mu$m$^2$; setup +0.0003~ns &
Fixed-T10 logic-only DC point; PT diagnostic failed setup/hold
($-0.001154/-0.022628$~ns); no admitted speedup/power/system result \\
Checkpoint & ep34 AEE 1.199514; activity 72.891G &
Final selection sealed; E2--E8 activation-dependent hardware rebind pending \\
Table A & 0 admitted rows &
Decoder-complete, memory-inclusive system row pending \\
\bottomrule
\end{tabular}
\end{center}%
}


\title{Capacity-Constrained Typed-Source Execution for Event-Based Spiking Optical Flow}
\author{Anonymous submission}
\date{}

\begin{document}
\maketitle

\begin{abstract}
Event-based spiking optical flow mixes binary spikes, signed analog neuron
outputs, repeated products, and persistent timestep state.  Operation sparsity
alone therefore does not guarantee speedup once finite SRAM capacity, parent
reads, partial-sum traffic, and service bandwidth are charged.  We present a
capacity-constrained execution architecture with three cooperating mechanisms:
single-port-aware finite product capture (C1), a typed signed K8 source fabric
sharing accumulator and endpoint state (C2), and exact fixed-T10 neuron service
(C3).  On the frozen ep34 ten-sample ledger, C1 exposes a
1.6945$\times$ ten-sample
same-ledger component opportunity below a 240-KiB capacity bound\modeltag;
its nine-macro mapped component meets 3-ns setup and hold, passes Formality, and
consumes 22.07~nJ in a separately bounded 64-row directed window.  Against an
equal-bandwidth K1$\times$8 baseline, C2 improves cycles by 1.0167$\times$ while
delivering 4.551$\times$ higher throughput per logic area and 77.66\% lower
logic area.  Both fresh C2 synthesis axes meet setup; the K8 axis independently
passes Formality and PrimeTime setup but retains a $-23.259$-ps fast-corner hold
violation.  K1$\times$8 independent timing, matched hold closure, and mapped
power remain open.  C3 closes prelayout
setup and hold in both synthesis and independent PrimeTime analysis, with 11,180
mapped equivalence points and zero failures.  The final decoder-complete,
memory-inclusive whole-network speedup and energy row remains \pending{} until
the ep34-bound replay completes; component ratios are not multiplied to infer it.
\end{abstract}

\section{Motivation and problem statement}
Event-driven optical flow exposes sparse work, but sparsity is not automatically
speedup: finite SRAM capacity, parent reads, partial-sum traffic, signed analog
activations, and timestep state can each become the bottleneck.  Our target is
therefore a bounded component architecture whose claims survive equal-resource
comparison rather than an unconstrained operation-count reduction.

\paragraph{Thesis.}
This paper studies how much sparse opportunity can be captured under a
240-KiB-class storage constraint and a shared signed source protocol.  The final
checkpoint identity is sealed as ep34; the thesis remains conditioned on an
ep34-bound decoder-complete replay, which is currently \pending.

\paragraph{Contributions.}
(1) C1 converts repeated-product opportunity into an executable capture schedule
under finite capacity, one-read/write parent storage, and dead-write suppression.
(2) C2 defines one protocol for binary and signed analog sources, amortizes
Acc24 plus endpoint state across up to eight sources, and adds lossless
token-set weight-row broadcast without reusing signed products; its physical
comparison is against equal service bandwidth, not a single narrow lane.
(3) C3 implements
exact fixed-T10 service with independently checked setup, hold, and mapped
equivalence.  A common replay contract separates accepted sources, component
cycles, memory traffic, and whole-network latency so the three component results
cannot be multiplied into a synthetic system ratio.

\placeholder{Figure 1: network/operator envelope and the C1--C3 ownership boundary.
Add only after the final checkpoint and decoder trace are bound.}

% -----------------------------------------------------------------------------
\clearpage
\section{Architecture and common execution contract}

\subsection{Signed sparse-source protocol}
The common protocol carries a source value, destination context, terminal state,
and validity.  C2 accepts at most eight typed signed sources while sharing the
Acc24 state and external endpoint.  The paper must distinguish protocol density,
accepted sources, executed cycles, and physical throughput per area.

\subsection{Finite-capacity memory boundary}
C1 is evaluated under an identical sub-240-KiB external capacity charge.  The
frozen ep34 ten-sample ledger replays 51.84 million source rows over four
bottleneck convolutions.  The current 1.694510$\times$ value is a CPU-cycle model
opportunity, not RTL, mapped-gate, decoder-complete, or system speedup.

The key distinction from an unbounded product-reuse model is the capture gap:
useful repeated products must survive parent-state port conflicts, finite
residency, completion ordering, and partial-sum ownership.  We consequently use
the same capacity, service width, and charge model for the product, bit-sparse,
and strongest zero-skipping baselines.  An official-artifact opportunity result
is reported separately from the executable C1 row.

\subsection{Neuron and decoder boundary}
C3 implements exact fixed-T10 service.  Decoder address-timed rows and the final
checkpoint rebind are not yet admitted into Table~A.  Consequently this skeleton
does not infer an Amdahl-weighted end-to-end ratio.

\placeholder{Figure 2: one source descriptor flowing through C2, C1 parent/product
capture, Acc24 commit, and C3 timestep service.  Annotate all SRAM ports and
ownership transitions.}

\paragraph{Invariant checklist.}
All final configurations must use one checkpoint identity, one resource model,
identical bandwidth and capacity, exact denominator definitions, and separately
reported model/RTL/pre-macro/mapped evidence.

% -----------------------------------------------------------------------------
\clearpage
\section{Component mechanisms}

\subsection{C1: finite-capacity product capture}
The mechanism combines repeated-product capture with single-port parent-state
scheduling and dead-write suppression.  Current CPU same-ledger cycles are
382,848,700 versus 648,741,051, or 1.694510$\times$\modeltag; the
same-coordinate bit-sparse denominator gives 1.688968$\times$.  A separate
28-nm, 3.0-ns, nine-macro component reports 166,514.312~$\mu$m$^2$,
PrimeTime setup/hold WNS of +0.027871/+0.001827~ns, and 16,549/16,549 passing
Formality points.  Public-port mapped SAIF/PTPX reports 29.0763~mW and
22.0689~nJ over a 64-row, 253-cycle directed window, including the nine SRAM
Liberty leaves.  This is a mixed-corner, prelayout, ideal-clock/ZeroWireload
component estimate without SPEF; it is not energy per frame or silicon power.
The CPU ratio, functional test, and physical/energy points remain separately
labelled.

\subsection{C2: shared typed signed K8 fabric}
Across five frozen directed workloads, K8 takes 1,913 cycles versus 1,945 cycles
for an equal-bandwidth K1$\times$8 baseline: 1.01672765$\times$.  The admitted
registered-fault-matched logic-only synthesis reports 130,822.775 versus
585,534.972~$\mu$m$^2$, implying 4.550657$\times$ throughput/mm$^2$ and
77.6576\% lower logic-cell area.  Minimum reported synthesis setup slacks are
+0.0018/+0.0014~ns for K8/K1$\times$8.  An independent K8-only
PrimeTime/Formality attempt was quarantined after hold violations before the
K1$\times$8 axis ran; none of its diagnostic values enter the admitted
evidence.  Matched independent PrimeTime/Formality remains open.  A one-shot
buffer-only hold repair was rejected because its 141,886.71~$\mu$m$^2$ area was
8.457\% above the admitted K8 logic baseline, exceeding the 5\% gate.  Therefore
hold closure is not claimed.  The cycle, area-efficiency,
and area-reduction values must appear together.  The area efficiency is not a
4.551$\times$ sparse speedup, and no production power result exists.

A lossless TSBG weight-delivery specialization is evaluated inside C2 rather
than as a fourth contribution.  An independently replayed ep34 40-sample CPU
premodel covers 960 FC1/FC2 pairs, 11,040 frames, and 44.64 million tokens.
Against an ordinary persistent same-capacity LRU4 row buffer with identical
eight-bank 128-B/cycle weight service, K8 compute, and Acc24 commit work, the
selected B4 point reduces serialized modeled cycles from 29,410,200,623 to
11,607,113,795 (2.533808$\times$) and weight traffic by 65.1978\%; the minimum
per-sequence ratio over four sequences is 2.517177$\times$.  Its maximum
incremental explicit state is 10,164~B.  These values are a same-I/O/cache CPU
premodel, not a same-area result, and the captured diagnostic INT8 codewords are
not hardware-quantization authority.  The exact M1880/M803 B4 component now
passes directed Synopsys VCS arithmetic, independently stalled/reordered bank
traffic, protocol attacks, and reset recovery.  In that 48-row directed test,
weight-bundle requests fall from 576 to 144 and scalar-bank requests from 4,608
to 1,152 (75\%); the result does not report an exact cycle ratio.  Matched DC
including all control/state, SRAM timing, and energy remain open, so the
2.533808$\times$ CPU premodel is not upgraded to RTL, same-area, or system
speedup.  B2 is retained only as a low-state ablation and B8 only as an
unphysicalized upper DSE point.

\subsection{C3: fixed-T10 exact neuron service}
The logic-only 28-nm, 3.0-ns component reports 63,756.125879~$\mu$m$^2$.
Design Compiler setup/hold WNS are +0.000300/+0.034585~ns; independent
PrimeTime setup/hold WNS are +0.000299/+0.030474~ns with no reported
violations.  Gate-to-gate Formality checks 11,180 compare points with zero
failures.  These are prelayout, zero-wireload, macro-free component results;
they do not admit throughput, speedup, power, energy, or a direct RTL-to-gate
equivalence claim.  Rank-adaptive paths are outside the present claim.

\placeholder{Figure 3: C1/C2/C3 microarchitecture panels.  Show finite storage,
equal-bandwidth K1$\times$8, and fixed-T10 denominators explicitly.}

% -----------------------------------------------------------------------------
\clearpage
\section{Evaluation plan and admitted evidence}

\subsection{Evidence registry}
\EvidenceRegistry

\subsection{Checkpoint and system status}
The final selected checkpoint is ep34, with AEE 1.199514 and 72.891G spike
activity.  Selection and the 40-sample ep34 activation capture are independently
sealed.  The latter contains 9,880 ordered records, 320 retained operator
payloads, 480 attention payloads, and all 93 live ATLIF modules, with zero
checkpoint-load mismatch.  This closes the activation-identity gate but does not
authorize a hardware speedup claim or retroactively rebind the older C1 cycle
ledger.  The production Table~A currently has zero
admitted rows; decoder-complete address-timed replay, memory traffic, total
latency, power, and energy remain \pending.

\subsection{Required final protocol}
The final evaluation must report: (i) a dense/bit-sparse/equal-resource baseline
ladder; (ii) decoder-complete cycles and bytes; (iii) component and whole-network
area, power and energy with explicit SRAM assumptions; (iv) at least three DSEC
sequences; and (v) AEE under the exact checkpoint used to generate activity.

\placeholder{Table A: final-checkpoint, decoder-complete, memory-inclusive system
rows.  Deliberately empty until a sealed replay supplies all denominators.}

% -----------------------------------------------------------------------------
\clearpage
\section{Comparison methodology and related work}

\subsection{Network-specific accelerator comparison}
Because the optical-flow network is co-designed with the accelerator, a direct
``same network on every published chip'' comparison is generally unavailable.
The final paper will separate three levels: our same-network baseline ladder;
iso-workload replay through public artifacts where supported; and normalized
published silicon/synthesis metrics with technology, precision, timestep,
        memory scope and workload caveats.

\paragraph{Metric definitions.}
Component speedup is baseline component cycles divided by proposed component
cycles on the same trace and resource point.  Throughput/mm$^2$ uses the same
accepted-work count and target clock for both implementations.  Whole-network
speedup is computed only from a decoder-complete address-timed replay; it is not
an Amdahl estimate assembled from unrelated component experiments.  Logic area,
macro area, dynamic power, SRAM energy, and off-chip energy remain separate
columns until a common integration row can sum them without double charging.

\subsection{Prosperity, Phi, and sparse SNN engines}
Prosperity~\cite{prosperity} combines unstructured spike processing, product
reuse, and low-overhead dispatch.  Its reported 7.4$\times$ average speedup is
against PTB over the paper's supported operator scope; the same-architecture
product-over-bit ablation is the more relevant comparison for C1.  Phi~\cite{phi}
uses a two-level pattern representation and reports a 3.45$\times$ system
speedup over Stellar.  Its 240-KiB storage point motivates our capacity-matched
evaluation, but its binary patterns differ from H67's typed signed sources.
FireFly-T~\cite{fireflyt} instead exposes fine-grained spike parallelism with a
multi-spike decoder and load-balanced execution; it motivates C2's equal-service
area comparison rather than a comparison against one low-bandwidth K1 lane.

These works illustrate why an operation reduction, a component service rate,
and a system speedup must be separated.  Any number obtained from a paper or
official artifact is therefore labeled as an original-work or
official-artifact result.  It is not called our RTL speedup, multiplied by
C1--C3, or inserted into the whole-network row without a matched
workload/resource denominator.  Error-constrained pruning such as
Bishop~\cite{bishop} is orthogonal to the present exact path; a future lossy row
would require a separately bound checkpoint and accuracy Pareto curve.

\begin{center}
\scriptsize
\begin{tabular}{p{0.20\columnwidth}p{0.31\columnwidth}p{0.39\columnwidth}}
\toprule
Comparison level & Allowed result & Mandatory qualifier \\
\midrule
Same network & C1--C3 and final Table A & Same checkpoint/resource/decoder scope \\
Official artifact & Artifact cycle opportunity & Original framework; not our RTL \\
Published paper & Reported throughput/energy & Technology, precision, memory, task \\
\bottomrule
\end{tabular}
\end{center}

\placeholder{Final comparison table: include Prosperity/Phi only with original
attribution and denominator footnotes; never imply a fabricated same-network run.}

% -----------------------------------------------------------------------------
\clearpage
\section{Limitations, closure checklist, and conclusion}

\subsection{Open admission gates}
\begin{enumerate}
  \item Preserve the sealed ep34 checkpoint, AEE, and 40-sample activation identity
        while rebinding each performance row to that exact trace.
  \item Complete decoder address-timed replay and populate the first Table-A row.
  \item Complete the K1$\times$8 Formality/PrimeTime axis and, if resources permit,
        one matched post-CTS/post-route K8/K1$\times$8 hold-repair experiment;
        do not repeat the rejected prelayout global buffer sweep.  Add C2 mapped
        power and matched memory timing.  C3 setup/hold is closed only at the
        prelayout logic-only boundary.
  \item TSBG-B4 directed VCS is closed.  Complete a matched ordinary-LRU4 versus
        broadcast-LRU4 DC/SRAM/energy comparison; until then, retain its
        2.5338$\times$ value only as an independently checked CPU premodel and
        its 75\% directed request reduction only as component work conservation.
  \item Produce component SAIF/PTPX and memory-inclusive energy without scaling an
        unrelated selected slice.
  \item Report a direct whole-network ratio only after all operators and stalls use
        the same resource model.
\end{enumerate}

\subsection{Claim firewall}
The abstract and conclusion may retain 1.6945$\times$ only with the CPU-model and
non-RTL qualifier; 4.551$\times$ only as equal-bandwidth throughput per logic area
alongside 1.0167$\times$ cycles and $-77.66\%$ logic area, with C2 hold explicitly
open; and all C3 numbers only
as prelayout logic-only area/timing/equivalence.  AEE 1.199514 and 72.891G are sealed ep34 selection facts, but they
do not authorize a hardware speedup claim merely because the ep34 activation
capture now passes independent audit; the same-resource replay must also close.
Table~A remains empty rather than extrapolated.

\subsection{Provisional conclusion}
The present evidence supports a component-architecture paper skeleton, not yet a
decoder-complete optical-flow accelerator claim.  The final conclusion will be
written only after the checkpoint, memory, decoder and energy gates close.

\begin{thebibliography}{9}\footnotesize
\bibitem{prosperity} C. Wei et al., ``Prosperity: Accelerating Spiking Neural
  Networks via Product Sparsity,'' in \emph{HPCA}, 2025, pp. 806--820.
\bibitem{phi} C. Wei et al., ``Phi: Leveraging Pattern-based Hierarchical
  Sparsity for High-Efficiency Spiking Neural Networks,'' in \emph{ISCA}, 2025,
  pp. 930--943.
\bibitem{fireflyt} T. Li et al., ``FireFly-T: High-Throughput Sparsity
  Exploitation for Spiking Transformer Acceleration with Dual-Engine Overlay
  Architecture,'' arXiv:2505.12771, 2025.
\bibitem{bishop} B. Xu et al., ``Bishop: Sparsified Bundling Spiking
  Transformers on Heterogeneous Cores with Error-constrained Pruning,'' in
  \emph{ISCA}, 2025.
\end{thebibliography}

\end{document}
